\documentclass[11pt]{article}
\usepackage{amsmath,amssymb,amsthm}
\usepackage[margin=1in]{geometry}
\usepackage{booktabs}
\newtheorem{theorem}{Theorem}
\newtheorem{lemma}{Lemma}
\newtheorem{corollary}{Corollary}
\theoremstyle{definition}\newtheorem{definition}{Definition}
\theoremstyle{remark}\newtheorem{remark}{Remark}

\title{Blackouts Preserving Rectangle Identifiability on Grid Points\\
\large A partial solution to Problem 001 of Kagey's Open Problem Collection}
\author{Arjun Pemmasani\thanks{Prepared with AI assistance (Claude, Anthropic);
all proofs re-derived by hand and all computations cross-verified by two
independent implementations. Code: [repo link].}}
\date{Preliminary draft --- \today}

\begin{document}
\maketitle

\begin{abstract}
[Stub. State: the reduction to a minimum hitting set of the difference hypergraph;
the Strip Theorem for $2\times m$; computed values for small general grids;
the pre-registered predictions confirmed by the solver.]
\end{abstract}

\section{Conventions and the presentation map}
% Five-layer stack: commitment game, presentation map, injectivity,
% pairwise characterization (Lemma 1), hitting-set form.

\section{The reduction}
% Lemma 1 with the element-chase proof; corollary: empty presentation subsumed.

\section{Preprocessing the difference hypergraph}
% Deduplication; superset-row deletion, stated in the correct containment direction.

\section{The Strip Theorem}
% Statement (m >= 3), classification of difference sets on strips (where m >= 3 is used),
% two pigeonholes + exactness line, witness characterization and count, m = 2 boundary.

\section{Computations}
% Tables emitted by the solver land in tables/ and are \input here.
% \input{tables/values.tex}

\section{Open directions}
% 3x3 and tilted rectangles; the random variant.

\end{document}
